% =============================================================
% Section 8. Discussion and Open Problems
% Template: paper/v0.1/notes/narrative.md §4 §8
% Structure: §8.1 contributions recap / §8.2 OP 4.5-b /
%            §8.3 OP 8b-rigor / §8.4 MRI engineering outlook /
%            §8.5 acknowledgments
% =============================================================

\noindent
We have constructed a categorical bridge between finite element
exterior calculus, cellular sheaf signal processing, and topological
signal processing, instantiated it in three previously disjoint
domains, and verified its main equality numerically to machine
precision. We now restate the seven contributions in the order they
were established, give the precise content of the two open problems
that remain, and outline the engineering verification on the MRI
side that lies outside this paper.

% -------------------------------------------------------------
\subsection{Summary of contributions}\label{subsec:disc-contrib}

\begin{enumerate}[leftmargin=*,label=(C\arabic*)]
\item The triangle equivalence (\cref{thm:triangle-equiv}) reads one
  scalar---$\sigma_{\min}(A)/\sigma_{\max}(A)$---through three
  schools: FEEC singular value bounds (\cref{subsec:tri-ab}), TopSP
  Riesz frame conditions (\cref{rem:three-readings}), and Robinson
  sheaf cohomology (\cref{subsec:tri-ac}). The proof produces a
  candidate sheaf bandwidth that is not yet categorically intrinsic
  (\cref{def:sheafbw-candidate}), and \cref{subsec:tri-narrowing}
  precisely narrows Robinson's open problem to the construction of an
  intrinsic version.
\item The categorically intrinsic candidate
  $\sheafBW(\mathscr{A}) := \sqrt{\gamma_+(\delta^0)/\Gamma_+(\delta^0)}$
  (\cref{def:sheafBW}) is unitarily invariant
  (\cref{prop:sheafBW-unitary}). The well-definedness of the
  spectral gap quantities relies on the closed Hilbert complex
  structure of \cref{thm:closed-hilbert-complex}, which uses the
  finite-dimensional restriction of \cref{def:finite-dim-sheaf}.
\item On the abstract sampling cell complex $X^A$
  (\cref{def:Xsamp}), the sampling sheaf $\samps$
  (\cref{def:sampling-sheaf}) reduces the intrinsic invariant to the
  ambient ratio $\sigma_{\min}(A)/\sigma_{\max}(A)$ \emph{strictly}
  (\cref{thm:strict-equality}) on the well-conditioned locus
  $\sigma_{\min}(A) > 0$. The double $0$-cell construction
  $\{\sigma_B, \sigma_\infty\}$ is what places $X^A$ into the
  Hansen--Ghrist regularity framework
  \cite[Def.~1, Cond.~4]{HansenGhrist2019};
  \cref{rem:rank-deficient-locus} delineates the rank-deficient
  boundary on which the strict equality fails.
\item Three previously disjoint domains
  (\cref{sec:three_instances}) fit a common categorical template:
  TopSP ($V_B^{\TopSP}$, \cref{subsec:6-topsp}), FEM--EM ($V_S^{\FEM}$,
  \cref{subsec:6-femem}), and Brillouin zone tetrahedron quadrature
  ($V_{\occ}^{\BZ}$, \cref{subsec:6-bz}). The three-instance
  comparison is recorded in \cref{tab:6-instances}.
\item The empirical proposition (\cref{empirical:8b}) formulates the
  $S$--$I$ Pareto trade-off as a directional empirical statement
  supported by two empirical channels and a heuristic skeleton
  (\cref{tab:6-evidence}); it is not promoted to theorem status, and
  it is not invoked in the proof of \cref{thm:strict-equality}---the
  categorical bridge is independent of it.
\item The numerical verification (\cref{sec:numerical}) records
  $\mathrm{reldiff} \le 1.64 \times 10^{-14}$ across the three toy
  instances at $\mathtt{seed} = 20260507$ (with Instances~1 and~3
  attaining $\le 8 \times 10^{-16}$), with explicit scope: the script
  verifies the SVD-based identity and the code path, not the
  per-instance physical evidence
  (\cref{subsec:7-honest-statement,tab:7-evidence}).
\item Two independent open problems remain, formulated below.
\end{enumerate}

% -------------------------------------------------------------
\subsection{Open Problem: non-trivial dagger on the full category}\label{subsec:disc-op-dagger}

\Cref{op:non-trivial-dagger} asks whether the unitary subcategory
groupoid of \cref{def:unitary-subcat} extends to a non-trivial
$\dagger$ structure on the full category $\CellSh^{\fin}_{\Hilb}(X)$.
The substantive content of this question is the tension between two
candidate lifts: dagger compactness in $\Hilb$
\cite[\S 2.3]{HeunenVicary2019} does not lift directly to cellular
sheaves because the restriction maps are not in general isometries,
while the dagger mono structure of
\cite[\S 3]{LahtinenStenvall2020} does lift but produces only the
unitary subgroupoid rather than a structure on the full category.
\Cref{rem:trivial-dagger} explains why the naive lift collapses; the
construction of a non-trivial extension is left open.
\Cref{op:non-trivial-dagger} is independent of all theorems of this
paper.

% -------------------------------------------------------------
\subsection{Open Problem: constructive Pareto proof}\label{subsec:disc-op-pareto}

\Cref{op:8b-rigor} asks for a constructive Pareto-optimality proof
upgrading the empirical proposition \cref{empirical:8b} to a
theorem. The proposed mode of attack proceeds in three steps:
compactness of the $V$-parameter space under suitable boundary
conditions; convex relaxation of the joint map
$V \mapsto (S(X^A_V), I(X^A_V))$, with the dual optimization
encoding the trade-off; and closing the duality gap via spectral
arguments paralleling \cite{Pilleboue2015}. Resolution of this
problem would replace \cref{empirical:8b} by a theorem and the
heuristic skeleton by a constructive proof. \Cref{op:8b-rigor} is
independent of the categorical bridge and of all theorems of this
paper.

% -------------------------------------------------------------
\subsection{Engineering outlook on the MRI side}\label{subsec:disc-mri-outlook}

The triangle equivalence (\cref{thm:triangle-equiv}) and the strict
equality (\cref{thm:strict-equality}) reduce, on the MRI sampling
problem, to a singular value ratio computable directly from the
acquisition matrix $A$. The project-internal numerical study
recorded in \cref{tab:6-evidence} examines $35$ sampling
configurations at $N = 200$; the multi-$N$ scan
($N \in \{200, 300, 1000\}$) shows the top three configurations
(Cartesian, random, sparkling) and the bottom two (spiral, radial)
separated by two to three orders of magnitude in the singular value
ratio across the entire range.
A larger-scale engineering verification---with $N \approx 10^4$
acquisition points, real phantom data, and one representative
baseline drawn from each of the four schools surveyed in
\cref{sec:related}---is the natural next step on the MRI side. That
engineering verification is outside the scope of this paper.

% -------------------------------------------------------------
\subsection*{Acknowledgments}

[Acknowledgments placeholder.]
